\section{Benchmark Design}
\label{sec:design}

\subsection{Domains and categories}

COT Bench v1 covers two production domains. \textbf{Banking} (account management,
transactions, fraud detection, compliance, loan processing) is structured,
tool-heavy, and has strict correctness requirements, which makes it well suited to
deterministic state verification. \textbf{Customer success} (CRM, ticketing, health
scoring, escalation, onboarding) has ambiguous goals where empathy and scope
judgment matter and multiple back-end systems must be coordinated, which is where
the judge panel earns its place.

Each domain has the same five scenario categories, chosen to probe distinct failure
modes rather than to tile a capability space:

\begin{center}
\begin{tabular}{ll}
\toprule
\textbf{Category} & \textbf{What it tests} \\
\midrule
Adaptive tool use & User needs shift mid-conversation; the agent must adapt \\
Scope management & User requests things outside the agent's capabilities \\
Empathetic resolution & Emotionally charged situations needing tools \emph{and} empathy \\
Extreme scenario recovery & Tool errors or unexpected data; the agent must recover \\
Adversarial input mitigation & Misleading, ambiguous, or manipulative inputs \\
\bottomrule
\end{tabular}
\end{center}

\subsection{Scenario schema}

Scenarios are JSON files under \texttt{data/scenarios/\{domain\}/} following the
v0.2 schema. Every scenario carries a synthetic \emph{persona} (traits, tone,
detail level, background, ranging from cooperative to adversarial), a list of
\emph{user goals} (3--10 interconnected goals, typically 5--8, designed so
completing one reveals information needed for another), a subset of domain
\emph{tools} with full JSON Schemas, and an in-character \emph{initial message} that
does not leak goal text or canonical IDs.

The v0.2 additions are what make a third of efficacy judge-independent. Each
scenario carries a \texttt{ground\_truth} block: a canonical world (balances,
records, customer state) the tool simulator must answer from and mutate. Keys are
canonical IDs (for example \texttt{BUS-CHK-001}) while the user-facing text stays
natural-language (``my business checking''), so resolving the friendly name to the
canonical ID is itself part of what the agent is graded on. Each scenario also
carries \texttt{expected\_state\_changes}, a list of deterministic assertions over
the post-conversation world expressed in a tiny vocabulary: \texttt{equals},
\texttt{increased\_by}, \texttt{decreased\_by}, and \texttt{contains} (a partial-dict
match over a list, with optional substring matching). For judgment calls the
benchmark grades the \emph{action}, not the \emph{outcome}: it asserts that a
fee-waiver request was \emph{submitted}, not that the waiver was approved, because
approval is the bank's decision and not the agent's. An empty assertion list is
meaningful: it encodes a no-unauthorized-mutation contract (the run passes if and
only if the world is unchanged), which is the right objective signal for the
refusal-heavy scope-management and adversarial categories.

\subsection{User and tool simulation}

Each scenario runs for up to ten user turns. Within a turn the loop is:
(1) a user simulator (GPT-4.1-mini, temperature 0.7) produces an in-character
message pursuing unmet goals; (2) the agent under test (temperature 0.0) responds,
with tools provided through native function calling (LangChain \texttt{bind\_tools})
rather than a regex-parsed JSON-in-text convention, so the benchmark measures real
tool-calling ability and does not penalize models tuned for function-calling APIs;
(3) if the agent calls tools, a tool simulator (GPT-4.1-mini, temperature 0.0)
returns schema-conforming responses and the agent is re-invoked on those results;
(4) steps 2--3 repeat within the turn until the agent emits a user-facing message
or an inner cap of five tool rounds is hit; (5) the user simulator decides whether
it is done talking. A content-embedded fallback parser exists only for providers
that return no native tool calls, and it logs a warning whenever it fires so its
use is measurable.

\paragraph{Stateful tool simulation.} For scenarios with a \texttt{ground\_truth},
the world is seeded from a deep copy of that block at the start of each run (one
world per run, reset between reliability repeats so mutations never leak across
runs). On every tool call the simulator is shown the current world and instructed
to answer only from it, returning a structured
\texttt{\{"response": \ldots, "state\_delta": \ldots\}}; mutating calls return a
dotted-path delta the runner applies deterministically, so a later ``check my
balance'' reflects an earlier transfer. Only the \texttt{response} is fed back to
the agent. This removes the simulator-incoherence failure mode in which asking a
balance twice yields two different numbers and the judge then penalizes the
\emph{agent} for the \emph{simulator's} drift.

\paragraph{Decoupling completion from goal-completion.} The published literature is
clear that simulated users are miscalibrated and over-cooperative, and that which
model plays the user can swing agent success by up to roughly nine
points~\cite{lostinsim2026,sim2real2026}. If the same simulator both plays the
customer and decides the goals are met, its bias contaminates the stop condition.
COT Bench therefore splits the two signals. The user simulator only declares that
it is \emph{done talking} (recorded as \texttt{ended\_by}). Whether the goals were
actually accomplished is the deterministic state check, evaluated against the
mutated world at the moment the conversation ended
(\texttt{state\_progress\_at\_end}). When the simulator ends a conversation while
the state check is still below 1.0, the run is flagged \texttt{premature\_end}: the
benchmark neither silently treats that ending as success nor extends the
conversation (which would let an already-biased simulator keep driving turns).
The discrepancy is made visible and aggregable, and the board reports a per-model
premature-ending rate. A planned part-two sim-sensitivity check will re-run a
subset under a second simulator and publish the score delta, directly quantifying
the swing the literature warns about.

\subsection{Authorship, contamination guards, and corpus statistics}

Every v0.2 scenario records an \texttt{authorship} block (\texttt{author\_model},
optional batch id and human reviewer). A hard guard refuses to generate any
scenario whose \texttt{author\_model} matches a model under test, with a
family-aware match so a different snapshot of a contestant still counts as that
contestant; a model under test never authors its own exam. The pipeline supports
multi-author generation across labs to avoid a single-author monoculture, runs a
cross-scenario dedup gate within each domain (near-duplicate initial-message and
goal-set overlap), and emits a distribution report over categories, difficulty,
persona reuse, and tool coverage. Because scenarios are synthetic, there is no
scraped public dataset that could have leaked into a model's training set; the
residual exposure (the public corpus is on GitHub) is addressed by the private
holdout described in Section~\ref{sec:governance}.

\paragraph{Corpus statistics (public split).} The public corpus is
92 scenarios, all on schema v0.2 with ground-truth worlds:
49 banking and 43 customer success. The category and difficulty
distribution is shown in Table~\ref{tab:corpus}. These counts are computed from
the committed corpus with \texttt{scripts/validate\_scenarios.py} (92/92
schema-valid) and should be regenerated the same way if the corpus changes before
the camera-ready version.

\begin{table}[h]
\centering
\caption{Public corpus distribution by category and difficulty. Values computed
from the committed \texttt{data/scenarios/} tree with
\texttt{scripts/validate\_scenarios.py}.}
\label{tab:corpus}
\begin{tabular}{lrr|lrrr}
\toprule
\textbf{Category} & \textbf{Bank} & \textbf{CS} & \textbf{Difficulty} & \multicolumn{3}{c}{\textbf{Count}} \\
\midrule
Adaptive tool use            & 12 & 10 & Easy   & \multicolumn{3}{c}{18} \\
Scope management             & 10 & 9  & Medium & \multicolumn{3}{c}{36} \\
Empathetic resolution        & 10 & 9  & Hard   & \multicolumn{3}{c}{38} \\
Extreme scenario recovery    & 8  & 7  &        & \multicolumn{3}{c}{} \\
Adversarial input mitigation & 9  & 8  &        & \multicolumn{3}{c}{} \\
\midrule
\textbf{Total}               & 49 & 43 & \textbf{Total} & \multicolumn{3}{c}{92} \\
\bottomrule
\end{tabular}
\end{table}

A separate private holdout (10 scenarios, authored to the same v0.2
schema and quality bar, stored outside the repository) is run alongside the public
set; only its aggregate gap is published (Section~\ref{sec:governance}).
